% !PRINTED 505
\weberpage{505}

Hieraus erhält man $\varkappa^2$ und $\varkappa'^2$ aus den Formeln [§ 54, (3)]
\[
\varkappa^2 = \frac{f_2(\omega)^8}{f(\omega)^8}, \quad \varkappa'^2 = \frac{f_1(\omega)^8}{f(\omega)^8},
\]
also
\[
\begin{aligned}
\omega &= \sqrt{-2}\colon & \varkappa^2 &= (\sqrt{2} - 1)^2, \\
\omega &= \sqrt{-3}\colon & \varkappa^2 &= \frac{2 - \sqrt{3}}{4} = \frac{(1 - \sqrt{3})^2}{8}, \\
\omega &= \sqrt{-7}\colon & \varkappa^2 &= \frac{8 - 3\sqrt{7}}{16} = \frac{(3 - \sqrt{7})^2}{32}.
\end{aligned}
\]

\begin{center}
§ 135.\qquad \gesperrt{Die Klasseninvariante $\varkappa^2$ und $\varkappa$.}
\end{center}

Es ist in § 126 gezeigt, daß, wenn die Diskriminante $D$ durch
$4$ teilbar ist, und die quadratische Gleichung, deren Wurzel $\omega$
ist, ungerade äußere Koeffizienten hat, $f(\omega)^{24}$ Klasseninvariante
für die Diskriminante $D$ ist.

Es ist aber nach § 54 und § 34
\weq{1}{\varkappa^2 \varkappa'^2 = \frac{16}{f(\omega)^{24}}, \quad \varkappa^2 = \frac{f_2(\omega)^8}{f(\omega)^8} = \frac{16}{f(\omega)^8\, f_1(2\omega)^8}.}
Genügt nun $\omega$ der Gleichung
\weq{2}{A\,\omega^2 + B\,\omega + C = 0}
mit der Diskriminante
\[
D = B^2 - 4\,A\,C,
\]
so genügt $\omega' = 2\,\omega$ der Gleichung
\[
A\,\omega'^2 + 2\,B\,\omega' + 4\,C = 0
\]
mit der Diskriminante $4\,D$, und nach § 126, (13) ist daher
$f_1(2\omega)^{24}$ Klasseninvariante für diese Diskriminante. Es ist also
$\varkappa^2 \varkappa'^2$ und $\varkappa^6$ rational ausdrückbar durch Klasseninvarianten der
Diskriminante $D$ und $4\,D$.

Die identische Gleichung
\weq{3}{\frac{\varkappa^2 \varkappa'^2 + \varkappa^4 \varkappa'^4 + 2\,\varkappa^6}{1 + \varkappa^6} = \varkappa^2}
zeigt, daß dasselbe auch für den Modul $\varkappa^2$ gilt.

Da $\varkappa^2$ durch jede lineare Substitution in eine rationale
(linear gebrochene) Funktion von $\varkappa^2$ übergeht, so kann jetzt auch
die Voraussetzung, daß die Koeffizienten $A$, $C$ ungerade sein
sollen, fallen gelassen werden, und wir haben den Satz:

2. \gesperrt{Ist die Diskriminante $D \equiv 0 \pmod{4}$, so ist $\varkappa^2(\omega)$ Klasseninvariante der Diskriminante $4\,D$.}
